\documentclass[10pt]{article}

\usepackage[utf8]{inputenc}

\usepackage[T1]{fontenc}

\usepackage{amsmath}

\usepackage{amsfonts}

\usepackage{amssymb}

\usepackage[version=4]{mhchem}

\usepackage{stmaryrd}

\usepackage{graphicx}

\usepackage[export]{adjustbox}

\graphicspath{ {./images/} }



\begin{document}

такой, что для любого $x \in G$ и любого $y \in G_{i}$ выполнено



\[

x y x^{-1} y^{-1} \in G_{i+1}

\]



здесь через $e$ обозначена единица группы $G$. Подгруппы $G_{i}$ могут быть выбраны нормальными. Любая Н.г. является разрешимой группой. Связная группа Ли нилыпотентна тогда и только тогда, когда нильпотентна ее алгебра Ли. Простейший пример нилыотентной группы Ли - группа всех верхнетреугольных матриц с единицами на главной диагонали. Описаны унитарные представления нильпотентных групп Ли. НИЛЬРАДИКАЛ конечномерной алгебры Лисм. Ли алгебра.



HOBИKOBA ГИПОтЕЗА в алгебраической геометрии - см. Алгебро-геометрические методы в теории поля.



HOBИKOBA-BECEЛOBA YPABHEHИE - уравнение



\[

\begin{aligned}

& V_{t}=\partial^{3} V+\bar{\partial}^{3} V+\partial(u V)+\bar{\partial}(\bar{u} V), \\

& \bar{\partial} u=3 \partial V, \quad \partial=\frac{\partial}{\partial z}, \quad z=x+i y ;

\end{aligned}

\]



было предложено в [1], как уравнение, описываюшее деформации потенциала $V(x, y, t)$ оператора Шредингера



\[

L=\partial \bar{\partial}+V

\]



сохраняющие спектральные данные оператора $L$ при одном уровне энергии. Такие деформации эквивалентны операторным уравнениям вида:



\[

\partial_{t} L=[L, A+\bar{A}]+(B+\bar{B}) L

\]



для Н.-В.у. операторы $A$ и $B$ равны: $A=\partial^{3}+u \partial, B=\partial u$. Н.-В. у. является двумерным аналогом Кортевега-де Фриса уравнения. Для него построены квазипериодич. конечнозонные решения, построены точные многосолитонные решения (см. [2]) и развита общая схема интегрирования задачи Коши для быстро убывающих начальных данных.



Лum.: [1] В е с ело в А. П., Но в и к в С.П., «Доклады АН СССР», 1984, т. 279, № 1, с. 20-24; [2] их же, там же, 1984, т. 279, № 4, с. 784-88.

И. М. Кричевер. HOBИКОВА-КРИЧЕВЕРА УРАВНЕНИЕ - уравнение, описывающее конечнозонные решения Кадомцева-Петвиашвили уравнения, отвечающие двумерным голоморфным расслоениям над эллиптической кривой. В алгебраич. форме имеет вид:



\[

\begin{gathered}

v_{t}=\frac{1}{4} v_{x x x}+\frac{3}{8 v_{x}}\left(v_{x x}^{2}-P_{3}(v)\right), \\

P_{3}(v)=4 v^{3}-g_{2} v-g_{3} .

\end{gathered}

\]



\section{Единственное уравнение вида}

\[

v_{t}=\frac{1}{4} v_{x x x}+f\left(v, v_{x}, v_{x x}\right)

\]



обладающее скрытой симметрией и не сводящееся к уравнению Кортевега-де Фриса преобразованиями типа Миуры.



Лum.: [1] К р и ч е в е р И. М., Н о в и к о в С. П., «Функщиональный анализ и его прилож.», 1978, т. 42, № 4, с. 41-52.А. П. Веселов. HOPMА в е к то р а - см. Евклидово пространство, Гильбертово пространство.



HOPMAIИ3ATOP подгр у п пы $H$ в группе $G$ - множество всех $g \in G$ таких, что для любого $h \in H$ выполнено соотношение $g h g^{-1} \in H$. Н. является подгруппой в $G$.



Ю. А. Неретин.



HOPMAЛЬНАЯ KPИВИЗНА - см. Главная кривизна.



НОРМАЛЬНАЯ ПОДГРУППА - см. Нормальный делитель. НОPМАЛЬНАЯ СЛУЧАЙНАЯ ФУНКЦИЯ - см. Гауссова случайная функция.



HOPMATbHAA DOPMA элли птич е ского ин тегр а л а - см. Эллиптический интеграл.



HOPMAJbHOE OTOSPAXEHИE - лагранжево отображение, сопоставляющее вектору $u v$ нормали к подмногообразию в евклидовом пространстве, приложенному в точке $u$, точку $v$ самого пространства. Каустика Н. о.- множество центров кривизны подмногообразия. Ограничение Н. о. на множество векторов длины $t$ является лежандровым, фронт к-рого называется эк в и д ист а нто й. При увеличении $t$ особенности эквидистант заметают каустику. М. Э. Казарян.



HOPMAЛЬHOE ПPOИЗВЕДEНИЕ О П раторов в к в а н т о в о й т е о р и и - запись произведения операторов в виде, когда все операторы рождения стоят слева от всех операторов уничтожения. Н. П. возникает в методе вторичного квантования; при этом предполагается, что любой оператор представим в виде полинома по операторам рождения и уничтожения. Отличительное свойство Н. п.- равенство нулю вакуумного среднего от любого оператора, записанного в виде Н. п. и не содержащего слагаемого, кратного единичному оператору. Н. п. было введено Дж. Виком (G. Wick, 1950) для того, чтобы исключить из квантовой теории поля формальные бесконечные величины типа энергии и заряда вакуумного состояния. Понятие Н. п. оказывается основным при решении многих фундаментальных вопросов квантовой теории поля таких, каК вывод фейнмановской диаграммной техники (см. Фейнмана диаграмма), установление связи между операторным формализмом и формализмом функционального интеграла, при построении аксиоматич. квантовой теории поля и т. П.



Н. п. операторов $A_{1}, \ldots, A_{n}$ обозначается



\[

: A_{1}, \ldots, A_{n}:

\]



Все свойства обычного произведения (линейность и т. д.) остаются и для Н. П., К-рое, кроме того, обладает свойством перестановочности операторов под знаком Н. п., при этом операторы, подчиняющиеся статистике Бозе-Эйнштейна, оказываются перестановочными, а подчиняющиеся статистике Ферми-Дирака - антиперестановочными.



Все динамич. величины, зависящие от операторов с одинаковыми аргументами (лагранжиан, тензор энергии-импульса, заряд и т. д.), во вторично-квантованной теории записываются в форме Н. П. Напр., оператор числа частиц для свободного скалярного поля $\varphi(x)$, удовлетворяющего Клейна-Гордона уравнению, в терминах операторов рождения $\varphi_{\mathbf{k}}^{+}$и уничтожения $\varphi_{\mathbf{k}}^{-}$частиц с импульсом $\mathbf{k}$ имеет вид:



\begin{center}

\includegraphics[max width=\textwidth]{2023_07_08_baca69decf6538bd240fg-1}

\end{center}



Для вакуумного среднего оператора $N$ получают



\[

(N)_{0}=\langle 0|N| 0\rangle=0

\]



так как $\varphi_{\mathbf{k}}^{-}|0\rangle=0$. Если бы $N$ не был представлен в виде Н. п., то выражение в скобках, возникающее из принципа соответствия с классич. теорией, привело бы к $(N)_{0}$, пропорциональному расходящемуся интегралу. Это типичный пример перестройки произведения в формализме Н. п. для операторов, подчиняюшихся статистике Бозе-Эйнштейна. В случае фермионов выражение в скобках имеет вид:



\[

\bar{a}_{\mathbf{k}, s}^{+} a_{\mathbf{k}, s}^{-}-\bar{a}_{\mathbf{k}, s}^{-} a_{\mathbf{k}, s}^{+}

\]



( $s$ - спиновая переменная), и для получения правильного оператора $N$, суммирующего все фермионные состояния, операторы рождения $\left(a^{+}\right)$и уничтожения $\left(a^{-}\right)$фермионов должны антикоммутировать под знаком Н. П. (черта над оператором означает дираковское сопряжение). Это - утверждение теоремы о связи спина и статистики (Паули теорема), вытекающей из принципа соответствия и формализма Н. п.



Для вычислений в квантовой теории поля необходимо установить связь Н. П. с обычным произведением и хронологииеским произведением. Эту связь устанавливают Вика теоремы. Спаривание двух линейных по операторам рождения и уничтожения операторов (соответственно хронологич. спаривание), обозначаемое ${\overrightarrow{A_{1}}}_{2}$, определяется как вакуумное среднее от обычного произведения (хронологич. произведения). Спаривание дается соответствующей пере-





\end{document}